\documentclass[11pt,a4paper]{article}

% ============================================================
% PREÁMBULO
% ============================================================

\usepackage[utf8]{inputenc}
\usepackage[T1]{fontenc}
\usepackage[spanish,es-nodecimaldot]{babel}

\usepackage{geometry}
\geometry{
    a4paper,
    margin=2.5cm
}

\usepackage{lmodern}
\usepackage{microtype}

\usepackage{xcolor}
\usepackage{listings}
\usepackage{booktabs}
\usepackage{longtable}
\usepackage{array}
\usepackage{amsmath}
\usepackage{amssymb}
\usepackage{hyperref}
\usepackage{enumitem}

% ============================================================
% CONFIGURACIÓN DE CÓDIGO
% ============================================================

\definecolor{codegray}{RGB}{245,245,245}
\definecolor{commentgreen}{RGB}{45,120,70}
\definecolor{keywordblue}{RGB}{40,70,150}

\lstdefinestyle{terminal}{
    backgroundcolor=\color{codegray},
    basicstyle=\ttfamily\small,
    breaklines=true,
    frame=single,
    columns=fullflexible
}

\lstdefinestyle{cpp}{
    language=C++,
    backgroundcolor=\color{codegray},
    basicstyle=\ttfamily\small,
    keywordstyle=\color{keywordblue},
    commentstyle=\color{commentgreen},
    stringstyle=\color{red!60!black},
    breaklines=true,
    frame=single,
    columns=fullflexible
}

\lstdefinestyle{textfile}{
    backgroundcolor=\color{codegray},
    basicstyle=\ttfamily\small,
    breaklines=true,
    frame=single,
    columns=fullflexible
}

\lstset{
    showstringspaces=false
}

% ============================================================
% HYPERREF
% ============================================================

\hypersetup{
    colorlinks=true,
    linkcolor=blue!60!black,
    urlcolor=blue!60!black,
    citecolor=blue!60!black,
    pdftitle={Priority Queue SKY130 -- Multi-Push Experimental Runbook},
    pdfauthor={},
}

% ============================================================
% DOCUMENTO
% ============================================================

\title{
    \textbf{Priority Queue SKY130}\\
    \large Multi-Push C++ Reference\\
    Resultados, Runbook y Glosario
}

\author{}
\date{\today}

\begin{document}

\maketitle

\tableofcontents
\newpage

% ============================================================
\section{Propósito del documento}
% ============================================================

Este documento registra la etapa experimental correspondiente a la
validación funcional de la referencia C++ de \emph{Multi-Push} para el
proyecto \texttt{priority-queue-sky130}.

El objetivo de esta etapa es comprobar que la implementación:

\begin{lstlisting}[style=textfile]
cpp/reference/multipush_reference.cpp
\end{lstlisting}

produce exactamente los mismos resultados observables que el
\emph{baseline} basado en un binary heap convencional.

La metodología adoptada es deliberadamente conservadora:

\[
\boxed{
\text{misma entrada}
\rightarrow
\text{misma salida}
\rightarrow
\text{verificación}
\rightarrow
\text{caracterización}
\rightarrow
\text{RTL}
}
\]

No se considera apropiado avanzar hacia una implementación RTL mientras
la referencia algorítmica no haya demostrado equivalencia funcional.

% ============================================================
\section{Contexto del proyecto}
% ============================================================

El proyecto busca estudiar una implementación de una \emph{hardware
priority queue} escalable y las optimizaciones asociadas al procesamiento
de múltiples operaciones \texttt{PUSH}.

La referencia principal del algoritmo es el trabajo:

\begin{quote}
\emph{A Fast Scalable Hardware Priority Queue and Optimizations for
Multi-Pushes}, de Samuel Collinson, Allan Bai y Oliver Sinnen,
Parallel and Reconfigurable Computing Lab, Department of Electrical,
Computer, and Software Engineering, University of Auckland, New Zealand.
\end{quote}

El proyecto utiliza una cadena progresiva:

\[
\boxed{
\text{C++ baseline}
\rightarrow
\text{C++ Multi-Push}
\rightarrow
\text{RTL}
\rightarrow
\text{síntesis}
\rightarrow
\text{SKY130}
}
\]

La etapa documentada aquí corresponde únicamente al segundo elemento:
la referencia C++ de Multi-Push.

% ============================================================
\section{Metodología de trabajo}
% ============================================================

El desarrollo y la ejecución experimental se realizan dentro del
entorno Docker utilizado por el proyecto.

Esto es importante para preservar la reproducibilidad del entorno de
EDA y evitar que diferencias entre el sistema anfitrión y el contenedor
afecten los resultados.

La separación conceptual es:

\begin{lstlisting}[style=textfile]
HOST
 |
 +-- Git / GitHub
 |
 +-- Docker
       |
       +-- C++ compilation
       +-- workload generation
       +-- baseline execution
       +-- Multi-Push execution
       +-- functional validation
       +-- characterization
       |
       +-- future RTL simulation
       +-- future synthesis
       +-- future SKY130 analysis
\end{lstlisting}

El repositorio Git permite versionar el proyecto, mientras que Docker
es el entorno de ejecución experimental.

% ============================================================
\section{Qué aprendimos}
% ============================================================

\subsection{La referencia Multi-Push ya puede ejecutar los workloads}

La implementación:

\begin{lstlisting}[style=cpp]
cpp/reference/multipush_reference.cpp
\end{lstlisting}

puede procesar los workloads generados para los tamaños:

\[
N \in \{8,16,32,64,128\}.
\]

Durante el experimento se han validado funcionalmente los cuatro primeros
tamaños:

\[
\boxed{N=8,16,32,64}
\]

El caso:

\[
N=128
\]

queda como siguiente punto de validación.

% ============================================================
\subsection{La equivalencia funcional se está comprobando contra el baseline}

Para cada tamaño se ejecutan dos modelos sobre exactamente el mismo
workload:

\[
\text{workload}
\rightarrow
\begin{cases}
\text{heap\_drain}\\
\text{multipush\_reference}
\end{cases}
\]

Posteriormente se comparan los archivos de salida.

La comprobación utilizada es:

\begin{lstlisting}[style=terminal]
cmp -s baseline.txt multipush.txt
\end{lstlisting}

Una salida como:

\begin{lstlisting}[style=textfile]
PASS: N=32 outputs identical
\end{lstlisting}

significa que los dos modelos produjeron exactamente la misma salida
observable para ese workload.

% ============================================================
\subsection{Resultados funcionales obtenidos}

Hasta este punto se obtuvo:

\begin{center}
\begin{tabular}{c c c c}
\toprule
$N$ & PUSH & POP & Equivalencia funcional \\
\midrule
8   & 8   & 8   & PASS \\
16  & 16  & 16  & PASS \\
32  & 32  & 32  & PASS \\
64  & 64  & 64  & PASS \\
128 & 128 & 128 & Pendiente \\
\bottomrule
\end{tabular}
\end{center}

Por lo tanto:

\[
\boxed{
N \in \{8,16,32,64\}
\Rightarrow
Output_{\mathrm{baseline}}
=
Output_{\mathrm{multipush}}
}
\]

Esto constituye evidencia de equivalencia funcional para esos workloads,
pero todavía no constituye una validación exhaustiva de todas las
posibles entradas.

% ============================================================
\subsection{Estadísticas observadas}

Las estadísticas obtenidas para Multi-Push fueron:

\begin{center}
\begin{tabular}{c r r r r r r}
\toprule
$N$ &
$C_{\mathrm{push}}$ &
$S_{\mathrm{push}}$ &
$C_{\mathrm{pop}}$ &
$S_{\mathrm{pop}}$ &
$C_{\mathrm{total}}$ &
$S_{\mathrm{total}}$ \\
\midrule
8  & 11  & 8  & 16  & 8   & 27  & 16 \\
16 & 21  & 10 & 56  & 28  & 77  & 38 \\
32 & 55  & 30 & 175 & 84  & 230 & 114 \\
64 & 116 & 59 & 452 & 223 & 568 & 282 \\
128 & -- & -- & -- & -- & -- & -- \\
\bottomrule
\end{tabular}
\end{center}

donde:

\begin{itemize}
    \item $C_{\mathrm{push}}$ representa comparaciones realizadas durante
    \texttt{HeapifyUp};
    \item $S_{\mathrm{push}}$ representa swaps realizados durante
    \texttt{HeapifyUp};
    \item $C_{\mathrm{pop}}$ representa comparaciones realizadas durante
    \texttt{HeapifyDown};
    \item $S_{\mathrm{pop}}$ representa swaps realizados durante
    \texttt{HeapifyDown}.
\end{itemize}

% ============================================================
\subsection{Para N=8}

El primer experimento produjo:

\begin{lstlisting}[style=textfile]
pushes=8
pops=8
push_comparisons=11
push_swaps=8
pop_comparisons=16
pop_swaps=8
total_comparisons=27
total_swaps=16
heapify_up_count=8
\end{lstlisting}

La comparación produjo:

\begin{lstlisting}[style=textfile]
PASS: N=8 outputs identical
\end{lstlisting}

% ============================================================
\subsection{Para N=16}

El experimento produjo:

\begin{lstlisting}[style=textfile]
pushes=16
pops=16
push_comparisons=21
push_swaps=10
pop_comparisons=56
pop_swaps=28
total_comparisons=77
total_swaps=38
heapify_up_count=16
\end{lstlisting}

Resultado:

\begin{lstlisting}[style=textfile]
PASS: N=16 outputs identical
\end{lstlisting}

% ============================================================
\subsection{Para N=32}

El experimento produjo:

\begin{lstlisting}[style=textfile]
pushes=32
pops=32
push_comparisons=55
push_swaps=30
pop_comparisons=175
pop_swaps=84
total_comparisons=230
total_swaps=114
heapify_up_count=32
\end{lstlisting}

Resultado:

\begin{lstlisting}[style=textfile]
PASS: N=32 outputs identical
\end{lstlisting}

% ============================================================
\subsection{Para N=64}

El experimento produjo:

\begin{lstlisting}[style=textfile]
pushes=64
pops=64
push_comparisons=116
push_swaps=59
pop_comparisons=452
pop_swaps=223
total_comparisons=568
total_swaps=282
heapify_up_count=64
\end{lstlisting}

Resultado:

\begin{lstlisting}[style=textfile]
PASS: N=64 outputs identical
\end{lstlisting}

% ============================================================
\section{Interpretación de los resultados}
% ============================================================

Los resultados permiten establecer una conclusión concreta:

\[
\boxed{
\text{La referencia Multi-Push produce la misma salida que el baseline
para } N=8,16,32,64.
}
\]

Sin embargo, todavía no corresponde concluir que Multi-Push presenta una
mejora algorítmica respecto del baseline.

La razón es que:

\[
\text{equivalencia funcional}
\neq
\text{mejora de coste}
\]

El siguiente experimento debe comparar las estadísticas internas del
baseline y de Multi-Push bajo exactamente los mismos workloads.

La pregunta pendiente es:

\[
\boxed{
\text{¿Dónde está realmente la reducción u organización diferente
del trabajo producida por Multi-Push?}
}
\]

Esta pregunta es especialmente importante antes de traducir el algoritmo
a RTL.

% ============================================================
\section{Interpretación de heap\_size y push\_location}
% ============================================================

La referencia Multi-Push utiliza dos conceptos separados.

\subsection{\texttt{heap\_size\_}}

Representa la cantidad de elementos que ya forman parte del heap
ordenado.

\[
\texttt{heap\_size\_}
=
\text{cantidad de elementos incorporados al heap}
\]

\subsection{\texttt{push\_location\_}}

Representa la posición donde se acepta el siguiente PUSH.

\[
\texttt{push\_location\_}
=
\text{primera posición libre para aceptar un PUSH}
\]

Por lo tanto:

\[
\boxed{
\texttt{push\_location\_} > \texttt{heap\_size\_}
}
\]

indica que existen elementos aceptados mediante PUSH que todavía no han
sido incorporados mediante \texttt{HeapifyUp}.

Esta separación es una de las partes conceptualmente importantes del
modelo.

% ============================================================
\section{Runbook}
% ============================================================

\subsection{Objetivo}

Este runbook describe cómo repetir la validación de la referencia
Multi-Push dentro del Docker.

La secuencia general es:

\[
\boxed{
\text{workload}
\rightarrow
\text{baseline}
\rightarrow
\text{Multi-Push}
\rightarrow
\text{compare}
\rightarrow
\text{statistics}
}
\]

% ============================================================
\subsection{Paso 1: entrar al entorno Docker}

El proyecto debe estar montado dentro del contenedor.

Entrar al proyecto:

\begin{lstlisting}[style=terminal]
cd /foss/designs/priority-queue-sky130
\end{lstlisting}

Verificar:

\begin{lstlisting}[style=terminal]
pwd
\end{lstlisting}

Resultado esperado:

\begin{lstlisting}[style=textfile]
/foss/designs/priority-queue-sky130
\end{lstlisting}

Entrar al directorio C++:

\begin{lstlisting}[style=terminal]
cd cpp
\end{lstlisting}

% ============================================================
\subsection{Paso 2: verificar los workloads}

Los workloads principales utilizan:

\[
\boxed{\mathrm{seed}=12345}
\]

y:

\[
\boxed{
N\in\{8,16,32,64,128\}
}
\]

Verificar:

\begin{lstlisting}[style=terminal]
ls -l workloads/generated/
\end{lstlisting}

Deben existir archivos como:

\begin{lstlisting}[style=textfile]
workload_8_seed12345.txt
workload_16_seed12345.txt
workload_32_seed12345.txt
workload_64_seed12345.txt
workload_128_seed12345.txt
\end{lstlisting}

% ============================================================
\subsection{Paso 3: ejecutar Multi-Push}

Para $N=8$:

\begin{lstlisting}[style=terminal]
./reference/multipush_reference \
    workloads/generated/workload_8_seed12345.txt \
    > /tmp/multipush_8.txt \
    2> /tmp/multipush_8.stats
\end{lstlisting}

Para $N=16$:

\begin{lstlisting}[style=terminal]
./reference/multipush_reference \
    workloads/generated/workload_16_seed12345.txt \
    > /tmp/multipush_16.txt \
    2> /tmp/multipush_16.stats
\end{lstlisting}

El mismo patrón se utiliza para los demás tamaños.

% ============================================================
\subsection{Paso 4: ejecutar el baseline}

Por ejemplo, para $N=16$:

\begin{lstlisting}[style=terminal]
./reference/heap_drain \
    workloads/generated/workload_16_seed12345.txt \
    > /tmp/baseline_16.txt
\end{lstlisting}

% ============================================================
\subsection{Paso 5: comparar outputs}

Comparar:

\begin{lstlisting}[style=terminal]
cmp -s /tmp/baseline_16.txt /tmp/multipush_16.txt \
    && echo "PASS: N=16 outputs identical" \
    || echo "FAIL: N=16 outputs differ"
\end{lstlisting}

También puede utilizarse:

\begin{lstlisting}[style=terminal]
diff -u /tmp/baseline_16.txt /tmp/multipush_16.txt
\end{lstlisting}

Si \texttt{diff} no produce salida, los archivos son idénticos.

% ============================================================
\subsection{Paso 6: inspeccionar estadísticas}

Por ejemplo:

\begin{lstlisting}[style=terminal]
cat /tmp/multipush_16.stats
\end{lstlisting}

Las estadísticas se mantienen separadas del output funcional:

\[
\boxed{
stdout=\text{resultado funcional}
}
\]

\[
\boxed{
stderr=\text{estadísticas}
}
\]

Esto permite comparar directamente los outputs sin contaminar la salida
con información de caracterización.

% ============================================================
\subsection{Paso 7: repetir para todos los tamaños}

El procedimiento completo es:

\begin{lstlisting}[style=terminal]
for N in 8 16 32 64 128; do

    ./reference/multipush_reference \
        "workloads/generated/workload_${N}_seed12345.txt" \
        > "/tmp/multipush_${N}.txt" \
        2> "/tmp/multipush_${N}.stats"

    ./reference/heap_drain \
        "workloads/generated/workload_${N}_seed12345.txt" \
        > "/tmp/baseline_${N}.txt"

    cmp -s \
        "/tmp/baseline_${N}.txt" \
        "/tmp/multipush_${N}.txt" \
        && echo "PASS: N=${N} outputs identical" \
        || echo "FAIL: N=${N} outputs differ"

done
\end{lstlisting}

% ============================================================
\subsection{Paso 8: revisar estadísticas}

Para revisar todos los resultados:

\begin{lstlisting}[style=terminal]
for N in 8 16 32 64 128; do
    echo "===== N=${N} ====="
    cat "/tmp/multipush_${N}.stats"
done
\end{lstlisting}

% ============================================================
\subsection{Paso 9: si aparece FAIL}

Si cualquier tamaño produce:

\begin{lstlisting}[style=textfile]
FAIL: N=XX outputs differ
\end{lstlisting}

el procedimiento debe detenerse.

No continuar todavía hacia:

\begin{lstlisting}[style=textfile]
RTL
synthesis
timing
area
power
SKY130
\end{lstlisting}

Primero comparar:

\begin{lstlisting}[style=terminal]
diff -u \
    /tmp/baseline_XX.txt \
    /tmp/multipush_XX.txt
\end{lstlisting}

y buscar el primer punto de divergencia.

% ============================================================
\section{Criterios de aceptación}
% ============================================================

La etapa C++ Multi-Push se considera funcionalmente validada cuando:

\begin{lstlisting}[style=textfile]
[PASS] mismo workload
[PASS] baseline ejecuta correctamente
[PASS] Multi-Push ejecuta correctamente
[PASS] N=8 outputs identical
[PASS] N=16 outputs identical
[PASS] N=32 outputs identical
[PASS] N=64 outputs identical
[PASS] N=128 outputs identical
\end{lstlisting}

Actualmente el estado es:

\begin{lstlisting}[style=textfile]
[PASS] N=8
[PASS] N=16
[PASS] N=32
[PASS] N=64
[PENDING] N=128
\end{lstlisting}

% ============================================================
\section{Qué debemos hacer después}
% ============================================================

Una vez validado $N=128$, el siguiente paso no será inmediatamente RTL.

Primero se realizará una caracterización comparativa:

\[
\begin{array}{c|c|c}
\text{Métrica} & \text{Baseline} & \text{Multi-Push}\\
\hline
\text{Push comparisons} & ? & ?\\
\text{Push swaps} & ? & ?\\
\text{Pop comparisons} & ? & ?\\
\text{Pop swaps} & ? & ?\\
\text{Total comparisons} & ? & ?\\
\text{Total swaps} & ? & ?
\end{array}
\]

El objetivo es determinar si la referencia implementada representa una
organización de trabajo distinta y útil para una posterior arquitectura
hardware.

En particular se deberá estudiar:

\begin{lstlisting}[style=textfile]
push_location_
heap_size_
busy()
finish_pending_pushes()
heapify_up()
heapify_down()
\end{lstlisting}

La transición posterior será:

\[
\boxed{
\text{C++ Multi-Push validado}
\rightarrow
\text{arquitectura RTL}
}
\]

% ============================================================
\section{Reglas metodológicas}
% ============================================================

Durante esta fase se mantienen las siguientes reglas:

\begin{enumerate}
    \item No cambiar el workload entre baseline y Multi-Push.
    \item Mantener la semilla fija.
    \item Ejecutar los experimentos dentro del Docker.
    \item Comparar primero outputs.
    \item Analizar estadísticas solamente después de establecer
    equivalencia funcional.
    \item No interpretar una estadística aislada como una mejora de
    hardware.
    \item No realizar afirmaciones de área, timing, potencia o PPA
    basadas únicamente en el modelo C++.
    \item No comenzar RTL hasta que el comportamiento algorítmico esté
    suficientemente comprendido.
\end{enumerate}

La regla de oro continúa siendo:

\[
\boxed{
\text{CORRECTNESS}
\rightarrow
\text{ALGORITHM}
\rightarrow
\text{HARDWARE}
}
\]

% ============================================================
\section{Glosario}
% ============================================================

\begin{longtable}{p{0.25\textwidth} p{0.68\textwidth}}
\toprule
\textbf{Término} & \textbf{Significado} \\
\midrule
\endhead

Priority Queue &
Estructura de datos que permite insertar elementos asociados a una
prioridad y extraer el elemento de mayor prioridad según el criterio
definido. \\

Binary Heap &
Representación basada en un árbol binario completo utilizada para
implementar eficientemente una priority queue. \\

HeapifyUp &
Operación que mueve un elemento hacia la raíz mientras tenga una
prioridad menor que la de su padre. \\

HeapifyDown &
Operación utilizada después de extraer la raíz para restaurar la
propiedad del heap desplazando un elemento hacia abajo. \\

PUSH &
Operación de inserción de un nuevo elemento en la priority queue. \\

POP &
Operación de extracción del elemento ubicado en la raíz del heap. \\

Multi-Push &
Optimización orientada a procesar múltiples operaciones PUSH con una
organización diferente del trabajo de mantenimiento del heap. \\

Baseline &
Implementación convencional utilizada como referencia funcional y
algorítmica. \\

Multi-Push Reference &
Modelo C++ que representa algorítmicamente la estrategia Multi-Push antes
de su traducción a RTL. \\

\texttt{heap\_size\_} &
Número de elementos que actualmente han sido incorporados al heap
ordenado. \\

\texttt{push\_location\_} &
Posición donde se almacenará el siguiente elemento aceptado por PUSH. \\

Pending PUSH &
Elemento que ya fue aceptado y almacenado pero todavía no ha sido
incorporado al ordenamiento del heap mediante HeapifyUp. \\

\texttt{busy()} &
Predicado que indica si existen PUSH pendientes. Es equivalente a
comprobar si \texttt{push\_location\_ > heap\_size\_}. \\

Functional Equivalence &
Condición en la cual dos implementaciones producen exactamente la misma
salida observable para una misma entrada. \\

Workload &
Archivo que contiene una secuencia de operaciones PUSH y POP utilizada
como entrada experimental. \\

Seed &
Semilla utilizada para generar determinísticamente un workload. En esta
caracterización se utiliza \texttt{12345}. \\

Comparison &
Operación de comparación entre dos elementos del heap. \\

Swap &
Intercambio de dos elementos almacenados en el vector del heap. \\

Characterization &
Medición del comportamiento interno del algoritmo, por ejemplo cantidad
de comparaciones y swaps. \\

RTL &
\emph{Register Transfer Level}; nivel de descripción utilizado para
diseñar hardware digital mediante Verilog o SystemVerilog. \\

SKY130 &
Proceso tecnológico open-source de 130 nm utilizado posteriormente para
síntesis y flujo físico. \\

Docker &
Entorno de ejecución utilizado para mantener controladas las
herramientas y dependencias del experimento. \\

Reproducibility &
Capacidad de repetir un experimento con las mismas entradas y entorno y
obtener resultados comparables. \\

\bottomrule
\end{longtable}

% ============================================================
\section{Glosario de comandos}
% ============================================================

\begin{longtable}{p{0.28\textwidth} p{0.63\textwidth}}
\toprule
\textbf{Comando} & \textbf{Función} \\
\midrule
\endhead

\texttt{cd} &
Cambia el directorio actual. \\

\texttt{pwd} &
Muestra el directorio actual. \\

\texttt{ls -l} &
Lista archivos mostrando información detallada. \\

\texttt{cat} &
Muestra el contenido de un archivo por la salida estándar. \\

\texttt{head} &
Muestra las primeras líneas de un archivo. \\

\texttt{wc -l} &
Cuenta las líneas de un archivo. \\

\texttt{./program} &
Ejecuta un programa ubicado en el directorio actual. \\

\texttt{>} &
Redirige \texttt{stdout} hacia un archivo. \\

\texttt{2>} &
Redirige \texttt{stderr} hacia un archivo. \\

\texttt{cmp} &
Compara dos archivos byte por byte. \\

\texttt{diff -u} &
Compara dos archivos mostrando las diferencias en formato unificado. \\

\texttt{for} &
Permite repetir comandos para múltiples tamaños de workload. \\

\texttt{g++} &
Compilador C++ utilizado para generar los ejecutables de referencia. \\

\texttt{git status} &
Muestra el estado del repositorio Git. \\

\texttt{git add} &
Agrega archivos al área de staging. \\

\texttt{git commit} &
Crea un commit con los cambios preparados. \\

\texttt{git push} &
Envía commits desde el repositorio local hacia el repositorio remoto. \\

\bottomrule
\end{longtable}

% ============================================================
\section{Estado experimental}
% ============================================================

El estado de esta etapa queda registrado como:

\begin{lstlisting}[style=textfile]
Project:
    priority-queue-sky130

Environment:
    Docker

Language:
    C++

Workload seed:
    12345

Validated sizes:
    N = 8
    N = 16
    N = 32
    N = 64

Pending:
    N = 128

Functional equivalence:
    PASS for N = 8,16,32,64

Algorithmic characterization:
    IN PROGRESS

RTL:
    NOT STARTED

Synthesis:
    NOT STARTED

SKY130 physical implementation:
    NOT STARTED
\end{lstlisting}

% ============================================================
\section{Conclusión}
% ============================================================

La etapa realizada demuestra que la referencia C++ de Multi-Push puede
procesar los workloads utilizados y que, para:

\[
N=8,16,32,64,
\]

produce exactamente los mismos outputs que el baseline.

Esto constituye el primer requisito para continuar con el proyecto:

\[
\boxed{
\text{equivalencia funcional}
}
\]

El siguiente objetivo experimental es completar:

\[
N=128
\]

y posteriormente realizar una comparación sistemática de las
operaciones internas entre baseline y Multi-Push.

Solamente después de comprender esa diferencia algorítmica se deberá
definir la arquitectura RTL.

Por lo tanto, el flujo metodológico permanece:

\[
\boxed{
\text{CORRECTNESS}
\rightarrow
\text{ALGORITHM}
\rightarrow
\text{HARDWARE}
}
\]

\end{document}
